\section{半小映射的分解定理（Xueqing Wen）}%
\label{sec:Decomposition theorem for semi-small maps (Xueqing Wen)}

\begin{quest}
    设 \(X\) 为（流形、簇、解析空间等）。我们如何计算 \(H^*(X, \C)\)？
\end{quest}

遗憾的是，这个问题太难了，所以我们把它换成一个新问题。

\begin{quest}
    我们如何用更小的空间来探测 \(H^*(X, \C)\)？
\end{quest}

\begin{thm}[K\"unneth 公式]
    我们有 \(H^*(X \times Y, \C) = H^*(X, \C) \otimes H^*(Y, \C)\)。
\end{thm}

遗憾的是，大多数空间并不是乘积，所以我们考虑下面的问题。

\begin{quest}
    设 \(f \colon X \to Y\) 为态射。我们能否用 \(Y\) 的上同调以及 \(f\) 的纤维的上同调来探测 \(X\) 的上同调？
\end{quest}

这个问题仍然非常困难，但至少我们可以研究它。我们从 \(f\) 是紧流形之间的纤维丛的情形开始。答案当然是 Leray 或 Serre 谱序列，但在到达那里之前，我们需要一些预备知识。

\begin{defn}
    \(Y\) 上的一个 \(\ul{\C}_Y\)-局部系是指 \(Y\) 上一个局部常值的向量空间层 \(\mc{L}\)。
\end{defn}

例如，若 \(Y\) 连通，则局部系对应于 \(\pi_1(Y)\) 的表示。另外，注意若 \(Y\) 是代数簇，我们需要在其解析化上工作。

\begin{exm}
    设 \(f\) 为如上所述的纤维丛，纤维为 \(F\)。则 \(R^i f_* \ul{\C}_X\) 是 \(Y\) 上以 \(H^i(F, \C)\) 为纤维的局部系。
\end{exm}

\begin{prop}
    我们有
    \[ H^*(X, \C) = H^*(X, \ul{\C}_X) = H^*(\mr{pt}, R a_{X*} \ul{\C}_X). \]
\end{prop}

\begin{thm}[Leray/Serre 谱序列]
    我们有谱序列
    \[ E_2^{p,q} = H^p(Y, R^q f_* \ul{\C}_X) \Rightarrow H^{p+q}(X, \C). \]
\end{thm}

我们现在考虑 \(f \colon X \to Y\) 是光滑射影簇之间的光滑射影映射的情形。

\begin{thm}[Deligne {\cite{Deligne1968}}]
    我们有非典范同构
    \[ R f_* \ul{\C}_X = \bigoplus_{q \geq 0} R^q f_* \ul{\C}[-q]. \]
    在上同调的层面上，我们有
    \[ H^*(X, \C) = \bigoplus_{q \geq 0} H^{*-q}(Y, R^q f_* \ul{\C}). \]
\end{thm}

这是射影簇的硬 Lefschetz 定理的一个推论，同时也表明 Serre 谱序列在 \(E_2\) 页退化。在拓扑中，这一结论对 Hopf 纤维化不成立。

我们现在考虑 \(f \colon X \to Y\) 是代数簇之间的本征映射的情形。此时我们有一个 \textit{Whitney 分层}\footnote{我们明确拒绝定义它。} \(Y = \bigsqcup_{\lambda} Y_{\lambda}\)，使得 \(f^{-1}(Y_{\lambda}) \to Y_{\lambda}\) 是可能奇异的簇之间的 \(C^{\infty}\) 纤维化。

\begin{defn}
    \(Y\) 上的 \(\ul{\C}_Y\)-模 \(\mc{F}\) 称为\textit{可构造层}，如果存在一个 Whitney 分层 \(Y = \bigsqcup Y_{\lambda}\)，使得对所有 \(\lambda\)，\(\mc{F}|_{Y_{\lambda}}\) 都是局部系。
\end{defn}

有了这一定义，我们就可以考虑 \(Y\) 上可构造层的导出范畴 \(\mc{D}_c^b(Y)\)。这一构造具有六函子形式体系，其中的函子为 \(f_*, f^*, f_!, f^!, \otimes, \mc{H}\mr{om}\)，并且通过函子
\[ D_Y \colon \mc{D}_c^b(Y) \to \mc{D}_c^b(Y). \]
满足 Verdier 对偶。

\begin{defn}
    复形 \(\mc{F}^{\bullet} \in \mc{D}_c^b(Y)\) 称为\textit{反常层}，如果它满足以下条件：
    \begin{enumerate}
        \item （支集条件）对所有 \(i\)，我们有 \(\dim \on{supp} \mc{H}^{-i}(\mc{F}^{\bullet}) \leq i\)；
        \item （余支集条件）对所有 \(i\)，我们有 \(\dim \on{supp} \mc{H}^{-i}(D_Y \mc{F}^{\bullet}) \leq i\)。
    \end{enumerate}
\end{defn}

\begin{rmks}\leavevmode
    \begin{enumerate}
        \item 反常层并不是层，而是次数位于 \([-\dim Y, 0]\) 中的复形。
        \item 它们有两个来源。一个是相交同调，另一个是 Riemann-Hilbert 对应。
    \end{enumerate}
\end{rmks}

我们将需要用到以下事实：\(Y\) 上反常层的范畴 \(\ms{Perv}(Y)\) 是一个 Abel 范畴，而且它还是 \(\mc{D}_c^b(Y)\) 上某个 \(t\)-结构的心。特别地，对任意 \(\mc{G} \in \mc{D}_c^b (Y)\)，我们可以考虑反常上同调 \(^p\mc{H}^i(\mc{G}) \in \ms{Perv}(Y)\)。

\begin{thm}[Beilinson-Bernstein-Deligne-Gabber {\cite{BBD1982}}]
    设 \(f \colon X \to Y\) 为代数簇之间的本征态射。假设 \(X\) 光滑且连通。\footnote{这并非必要，但若去掉光滑性假设，我们就需要考虑相交复形。}则
    \[ R f_* \ul{\C}_X [\dim X] = \bigoplus_{i} {^p\mc{H}^i}(R f_* \ul{\C}_X)[-i]. \]
    此外，每个 \(^p\mc{H}^i (R f_* \ul{\C}_X)\) 都是半单的（即每个子对象都是直和项）。
\end{thm}

\begin{defn}
    设 \(f \colon X \to Y\) 为本征满态射，其中 \(X\) 光滑且不可约。令 \(Y = \bigsqcup Y_{\lambda}\) 为一个分层，使得 \(f^{-1}(Y_{\lambda}) \to Y_{\lambda}\) 是 \(C^{\infty}\)-纤维化。我们称 \(f\) 是\textit{半小}的，如果对所有 \(y_{\lambda} \in Y_{\lambda}\)，都有 \(\on{codim}_Y Y_{\lambda} \geq 2 \dim f^{-1} (y_{\lambda})\)。若等号成立，则称 \(Y_{\lambda}\) 是\textit{相关}的。若唯一的相关层片是稠密的那个，则称 \(f\) 是\textit{小}的。
\end{defn}

\begin{thm}[{\cite{DeCataldoMigliorini2002}}]
    若 \(f \colon X \to Y\) 是半小的，则 \(R f_* \ul{\C}_X [\dim X]\) 是 \(Y\) 上的反常层。
\end{thm}

\begin{proof}
    对 \(y_{\lambda} \in Y_{\lambda} \subseteq Y\)，我们看到
    \[ R f_* \ul{\C}_X |_{y_{\lambda}} = \bigoplus H^i (f^{-1}(y_{\lambda}), \C). \]
    注意到当 \(i > 2 \dim f^{-1}(y_{\lambda}) \leq \on{codim}_Y (Y_{\lambda})\) 时它为零。对所有层片 \(Y_{\lambda}\) 变动，我们便得到支集条件。因为 \(f\) 是本征的，我们有 \(R f_* = f_!\)，再利用 \(D_Y R f_* = f_! D_X\)，便得到余支集条件。
\end{proof}

\begin{exms}
    我们现在给出一些半小映射的例子。
    \begin{enumerate}
        \item 曲面之间任何一般有限的本征映射都是半小的。
        \item 令 \(Y = \ab\{xy = t^{\ell+1}\} \subseteq \C^3\)。此时存在奇点的 crepant 解消 \(f \colon \tilde{Y} \to Y\)，且奇点上的纤维是一条 \(\P^1\) 的链，因此它是半小的（亦见前一个例子）。
    \end{enumerate}
\end{exms}

\begin{thm}[Kaledin {\cite{Kaledin2006}}]
    从非奇异全纯辛簇出发的射影双有理映射是半小的。
\end{thm}

\begin{exm}
    例如，Springer 解消与 Grothendieck-Springer 解消都是半小的。此外，曲面上点的 Hilbert 概形的 Hilbert-Chow 态射也是半小的。
\end{exm}

设 \(f \colon X \to Y\) 是半小的，并考虑一个分层 \(Y = \bigsqcup Y_{\lambda}\)。若 \(Y_{\lambda}\) 是相关的，考虑
\[ \mc{L}_{\lambda} = R^{d_{\lambda}} f_* \ul{\C}_{f^{-1}(Y_{\lambda})}, \]
其中 \(d_{\lambda} = 2 \dim f^{-1} (y_{\lambda})\)。则 \(\mc{L}_{\lambda}\) 是半单的。这里注意 \(\mc{L}_{|\lambda} = H^{d_{\lambda}}(f^{-1}(y_{\lambda}), \C)\)。现在令 \(\mr{IC}(\mc{L}_{\lambda})\) 为 \(\mc{L}_{\lambda}[\dim Y_{\lambda}]\) 在 \(\ol{Y}_{\lambda}\) 上的中间延拓，它是一个反常层。

\begin{thm}[{\cite{DeCataldoMigliorini2002}}]
    设 \(f \colon X \to Y = \bigsqcup Y_{\lambda}\) 是半小的。则有典范分解
    \[ R f_* \ul{\C}_X[\dim X] = \bigoplus_{Y_{\lambda}\text{ relevant}} \mr{IC}(\mc{L}_{\lambda}). \]
\end{thm}

\begin{exm}
    设 \(G\) 为约化群，其 Lie 代数为 \(\g\)，并取定一个 Borel 子群 \(B\)。则
    \[ G \times_B \b = \tilde{\g} \xrightarrow{\pi} \g \]
    是小的。因此，我们有
    \[ R \pi_* \ul{\C}_{\tilde{\g}} [\dim \g] = \mr{IC}(\mc{L}). \]
    回忆 \(\tilde{\g}^{\on{rs}} \to \g^{\on{rs}}\) 是一个 \(W\)-torsor。则对 \(a \in \g^{\on{rs}}\)，我们有 \(\mc{L}_{|a} \cong \C[W]\)，因此 \(\mc{L} = R^0 \pi_* \ul{\C}_{\tilde{\g}^{\on{rs}}}\)。这意味着 \(W\) 作用于 \(\mc{L}\)，而这一作用可以延拓到 \(\mr{IC}(\mc{L})\) 上，也可以延拓到 \(R \pi_* \C_{\tilde{\g}}\) 上，特别地，它作用于任意 \(a \in \g\) 处的纤维 \(R \pi_* \C_{\tilde{\g}}|_a\)。
\end{exm}

\begin{exm}
    考虑 Springer 解消 \(\mu \colon T^*(\on{GL}_3/B) \to \mc{N} = \mathbb{O}_{(1^3)} \sqcup \mathbb{O}_{(2,1)} \sqcup \mathbb{O}_{(3)}\)。该映射是半小的，但遗憾的是所有幂零轨道都是相关的。注意到 \((2,1)\) 轨道上的纤维是两个 \(\P^1\) 在一个结点处粘合的并，\((3)\) 轨道上的纤维是一个点，而 \((1^3)\) 轨道上的纤维是 \(\on{GL}_3/B\)。

    我们看到 \(\mc{L}_{(3)}\) 必为平凡的，因为 \(\mu\) 在该轨道上方是同构。然后，\(\mc{L}_{(2,1)} = R^2 \mu_* \ul{\C}\) 是 \(\mathbb{O}_{(2,1)}\) 上秩为 \(2\) 的局部系。最后，\(\mc{L}_{(1^3)} = R^6 \mu_* \ul{\C}\) 是 \(\mathbb{O}_{(1^3)} = \mr{pt}\) 上秩为 \(1\) 的局部系。因此，我们有
    \[ R \mu_* \ul{\C}_{T^*\on{GL}_3/B}[6] = \mr{IC}(\mc{L}_{(3)}) \oplus \mr{IC}(\mc{L}_{(2,1)}) \oplus \mr{IC}(\mc{L}_{(1^3)}). \]
    限制到 \(b \in \mathbb{O}_{(3)}\)，我们得到 \(\mc{L}_{(3)} = H^0(\mr{pt})\)。限制到 \(a \in \mathbb{O}_{(2,1)}\)，我们得到
    \[ H^0(X, \C) \oplus H^2(X, \C) = \mr{IC}(\mc{L}_{(3)})|_{(a)} \oplus \mr{IC}(\mc{L}_{(2,1)})|_{(a)}. \]
    限制到 \(0 = \mathbb{O}_{(1^3)}\)，我们得到
    \begin{align*}
        H^*(\on{GL}_3/B, \C) = \C[W] &= \mr{IC}(\mc{L}_{(3)})|_0 \oplus \mr{IC}(\mc{L}_{(2,1)})|_0 \oplus \mc{L}_{(1^3)}|_0 \\
        &= H^0(\on{GL}_3/B, \C) \oplus \ab(\text{一个 } 4 \text{ 维的直和项}) \oplus H^6(\on{GL}_3/B, \C).
    \end{align*}
    作为 \(W = S_3\) 的表示，它们分别对应于平凡表示、两份标准表示以及符号表示。关于 Springer 理论的构造以及这些分解，见~\cite{BorhoMacPherson1983}。
\end{exm}
